%%%%%%%%%%%%%%%%%%%%%%%%%%%%%%%%%%
%\vspace{1 in}
%\draftnewpage
%\section{Prologue}
\chapter*{Prologue}
%%%%%%%%%%%%%%%%%%%%%%%%%%%%%%%%%%

\bigskip

\begin{quotation}
\emph{\dots Cela suffit pour faire comprendre que dans les cinq m\'emoires des Acta mathematica que j'ai consacr\'es \`a l'\'etude des transcendantes fuchsiennes et klein\'eennes, je n'ai fait qu'effleurer un sujet tr\`es vaste, qui fournira sans doute aux g\'eom\`etres l'occasion de nombreuses et importantes d\'ecouvertes.\Footnote{This is enough to make it apparent that in these five memoirs in Acta Mathematica which I have dedicated to the study of Fuschian and Kleinian transcendants, I have only skimmed the surface of a very broad subject, which will no doubt provide geometers with the opportunity for many important discoveries.}}\flushright{ -- H. Poincar\'e, Acta Mathematica, {\bf 5}, 1884, p. 278.}
\end{quotation}

\smallskip

The theory of discrete subgroups of real hyperbolic space has a long history. It was inaugurated by Poincar\'e, who developed the two-dimensional (Fuchsian) and three-dimensional (Kleinian) cases of this theory in a series of articles published between 1881 and 1884 that included numerous notes submitted to the C. R. Acad. Sci. Paris,
%This is the full name of the journal: Comptes Rendus Hebdomadaires des S\'eances de l'Acad\'emie des Sciences (C. R. Acad. Sci. Paris Sér. A-B)
a paper at Klein's request in Math. Annalen, and five memoirs commissioned by Mittag-Leffler for his then freshly-minted Acta Mathematica. One must also mention the complementary work of the German school that came before Poincar\'e and continued well after he had moved on to other areas, viz. that of Klein, Schottky, Schwarz, and Fricke. See \cite[Chapter 3]{Gray2} for a brief exposition of this fascinating history, and \cite{Gray1, Saint-Gervais} for more in-depth presentations of the mathematics involved.

We note that in finite dimensions, the theory of \emph{higher-dimensional Kleinian groups}, i.e., discrete isometry groups of the hyperbolic $d$-space $\H^d$ for $d \geq 4$, is markedly different from that in $\H^3$ and $\H^2$. For example, the Teichm\"uller theory used by the Ahlfors--Bers school (viz. Marden, Maskit, J\o rgensen, Sullivan, Thurston, etc.) to study three-dimensional Kleinian groups has no generalization to higher dimensions. Moreover, the recent resolution of the Ahlfors measure conjecture \cite{Agol,CalegariGabai} has more to do with three-dimensional topology than with analysis and dynamics. Indeed, the conjecture remains open in higher dimensions \cite[p. 526, last paragraph]{Kapovich}. Throughout the twentieth century, there are several instances of theorems proven for three-dimensional Kleinian groups whose proofs extended easily to $n$ dimensions (e.g. \cite{BeardonMaskit,Mostow}), but it seems that the theory of higher-dimensional Kleinian groups was not really considered a subject in its own right until around the 1990s. For more information on the theory of higher-dimensional Kleinian groups, see the survey article \cite{Kapovich}, which describes the state of the art up to the last decade, emphasizing connections with homological algebra.

But why stop at finite $n$? Dennis Sullivan, in his IH\'ES \emph{Seminar on Conformal and Hyperbolic Geometry} \cite{Sullivan_seminar} that ran during the late 1970s and early '80s, indicated a possibility of developing the theory of discrete groups acting by hyperbolic isometries on the open unit ball of a separable infinite-dimensional Hilbert space.\Footnote{This was the earliest instance of such a proposal that we could find in the literature, although (as pointed out to us by P. de la Harpe) infinite-dimensional hyperbolic spaces without groups acting on them had been discussed earlier \cite[\627]{Michal}, \cite{Michal2, Harpe_thesis}. It would be of interest to know whether such an idea may have been discussed prior to that.} Later in the early '90s, Misha Gromov observed the paucity of results regarding such actions in his seminal lectures \emph{Asymptotic Invariants of Infinite Groups} \cite{Gromov4} where he encouraged their investigation in memorable terms:
%We quote from `$6$A.III: Infinite dimensional symmetric spaces and buildings' in \cite{gro}:
``The spaces like this [infinite-dimensional symmetric spaces] \ldots ~look as cute and sexy to me as their finite dimensional siblings but they have been for years shamefully neglected by geometers and algebraists alike''.

Gromov's lament had not fallen to deaf ears, and the geometry and representation theory of infinite-dimensional hyperbolic space $\H^\infty$ and its isometry group have been studied in the last decade by a handful of mathematicians, see e.g. \cite{BIM, DelzantPy, MonodPy}. However, infinite-dimensional hyperbolic geometry has come into prominence most spectacularly through the recent resolution of a long-standing conjecture in algebraic geometry due to Enriques from the late nineteenth century. Cantat and Lamy \cite{CantatLamy} proved that the Cremona group (i.e. the group of birational transformations of the complex projective plane) has uncountably many non-isomorphic normal subgroups, thus disproving Enriques' conjecture. Key to their enterprise is the fact, due to Manin \cite{Manin}, that the Cremona group admits a faithful isometric action on a non-separable infinite-dimensional hyperbolic space, now known as the Picard--Manin space.

Our project was motivated by a desire to answer Gromov's plea by exposing a coherent general theory of groups acting isometrically on the infinite-dimensional hyperbolic space $\H^\infty$. In the process we came to realize that a more natural domain for our inquiries was the much larger setting of semigroups acting on Gromov hyperbolic metric spaces -- that way we could simultaneously answer our own questions about $\H^\infty$ and construct a theoretical framework for those who are interested in more exotic spaces such as the curve graph, arc graph, and arc complex \cite{HPW, MasurSchleimer, HilionHorbez} and the free splitting and free factor complexes \cite{HandelMosher, BestvinaFeighn, KapovichRafi, HilionHorbez}. These examples are particularly interesting as they extend the well-known dictionary \cite[p.375]{Bestvina_survey} between mapping class groups and the groups $\Out(\F_N)$. In another direction, a dictionary is emerging between mapping class groups and Cremona groups, see \cite{BlancCantat, DHKK}. We speculate that developing the Patterson--Sullivan theory in these three areas would be fruitful and may lead to new connections and analogies that have not surfaced till now.

%Why do we care about infinite-dimensional space? Early on in our investigations, we learned that a theorem of Edelstein \cite[Theorem 2.1]{Edelstein} (see also \6\ref{subsubsectionedelstein} below) corresponds with a parabolic isometry of $\H^\infty$ which exhibits fairly erratic behavior -- the orbit of $\zero$ tends to the parabolic fixed point along a subsequence, while infinitely often returning to a neighborhood of the origin. This behaviour has no analogue in finite dimensions. In general, infinite-dimensional space is a wellspring of outlandish examples, such as those constructed using the Burger--Iozzi--Monod representation theorem \cite[Theorem 1.1]{BIM} (see also Theorem \ref{theoremBIM} below), and Poincar\'e extension from $\Isom(\HH)$ (see Section \ref{sectionparabolic}). Examples from $\H^\infty$ often help us to formulate our theorems optimally; for example, a certain counterexample to the naive generalization of the Bishop--Jones theorem, namely Example \ref{exampleAutTBIM}, helped us formulate the definition of the modified Poincar\'e exponent (Definition \ref{definitionmodifiedexponent}). Similarly, Proposition \ref{propositioncounterexample} showed us that the assumption of divergence type would be necessary in order to prove the Patterson--Sullivan theorem without the assumption of compact type. \commariusz{I would either remove this paragraph or would add more to it. It the present form it is not really convincing. We can talk more about it on the phone or Skype.}

In a similar spirit, we believe there is a longer story for which this monograph lays the foundations. In general, infinite-dimensional space is a wellspring of outlandish examples and the wide range of new phenomena we have started to uncover has no analogue in finite dimensions. The geometry and analysis of such groups should pique the interests of specialists in probability, geometric group theory, and metric geometry. More speculatively, our work should interact with the ongoing and still nascent study of geometry, topology, and dynamics in a variety of infinite-dimensional spaces and groups, especially in scenarios with sufficient negative curvature. Here are three concrete settings that would be interesting to consider: the universal Teichm\"uller space, the group of volume-preserving diffeomorphisms of $\R^3$ or a $3$-torus, and the space of K\"ahler metrics/potentials on a closed complex manifold in a fixed cohomology class equipped with the Mabuchi--Semmes--Donaldson metric. We have been developing a few such themes. The study of thermodynamics (equilibrium states and Gibbs measures) on the boundaries of Gromov hyperbolic spaces will be investigated in future work \cite{DSU2}.  We speculate that the study of stochastic processes (random walks and Brownian motion) in such settings would be fruitful. Furthermore, it would be of interest to develop the theory of discrete isometric actions and limit sets in infinite-dimensional spaces of higher rank.
%We also hope to study stochastic processes (random walks and Brownian motion) in such settings, and to develop the theory of limit sets in spaces of higher rank.

\bigskip

{\bf Acknowledgements.} This monograph is dedicated to our colleague Bernd O. Stratmann, who passed away on the 8th of August, 2015. Various discussions with Bernd provided inspiration for this project and we remain grateful for his friendship. The authors thank D. P. Sullivan, D. Mumford, B. Farb, P. Pansu, F. Ledrappier, A. Wilkinson, K. Biswas, E. Breuillard, A. Karlsson, I. Assani, M. Lapidus, R. Guo, Z. Huang, I. Gekhtman, G. Tiozzo, P. Py, M. Davis, M. Roychowdury, M. Hochman, J. Tao, P. de la Harpe, T. Barthelme, J. P. Conze, and Y. Guivarc'h for their interest and encouragement, as well as for invitations to speak about our work at various venues. We are grateful to S. J. Patterson, J. Elstrodt, and \'E. Ghys for enlightening historical discussions on various themes relating to the history of Fuchsian and Kleinian groups and their study through the twentieth century, and to D. P. Sullivan and D. Mumford for suggesting work on diffeomorphism groups and the universal Teichm\"uller space. We thank X. Xie for pointing us to a solution to one of the problems in our Appendix \ref{appendixopenproblems}. The research of the first-named author was supported in part by 2014-2015 and 2016-2017 Faculty Research Grants from the University of Wisconsin--La Crosse. The research of the second-named author was supported in part by the EPSRC Programme Grant EP/J018260/1. The research of the third-named author was supported in part by the NSF grant DMS-1361677.

\endinput